% Theorem 1: Effective Dimensionality of Evaluation Suites
% STEREOLOGY Paper — Proof Document

\subsection{Setup and Notation}

Let $\mathcal{M} = \{c_1, \ldots, c_n\}$ denote a population of $n$ models, each possessing a \emph{capability profile} $c(c_i) \in \mathbb{R}^D$ for some unknown $D$. A benchmark suite $\Pi = \{\pi_1, \ldots, \pi_k\}$ maps each profile to an observable score vector:
\[
  \Pi(c_i) = (\pi_1(c_i), \ldots, \pi_k(c_i)) \in \mathbb{R}^k.
\]
We observe the \emph{score matrix} $S \in \mathbb{R}^{n \times k}$ with $S_{ij} = \pi_j(c_i)$.

\begin{definition}[Effective Dimensionality]
Let $\Sigma = \mathrm{Corr}(S)$ denote the $k \times k$ sample correlation matrix of the score matrix, with eigenvalues $\lambda_1 \geq \lambda_2 \geq \cdots \geq \lambda_k \geq 0$. The \emph{effective dimensionality} of the benchmark suite is the participation ratio:
\[
  d_{\mathrm{eff}} = \frac{\left(\sum_{i=1}^k \lambda_i\right)^2}{\sum_{i=1}^k \lambda_i^2}.
\]
\end{definition}

\noindent\textbf{Properties.} Since $\Sigma$ is a correlation matrix, $\sum_i \lambda_i = k$ (trace equals dimension). Thus $d_{\mathrm{eff}} = k^2 / \sum_i \lambda_i^2$. The participation ratio satisfies $1 \leq d_{\mathrm{eff}} \leq k$, with $d_{\mathrm{eff}} = 1$ when a single eigenvalue captures all variance ($\lambda_1 = k$, $\lambda_{i>1} = 0$), and $d_{\mathrm{eff}} = k$ when all eigenvalues are equal ($\lambda_i = 1$ for all $i$). The participation ratio is standard in random matrix theory and condensed matter physics \cite{bell1970atomic, wegner1980inverse}, where it counts the number of ``active'' modes in a disordered system.

\begin{theorem}[Effective Dimensionality and Variance Capture]
\label{thm:deff}
Let $\mathcal{C} \subseteq \mathbb{R}^D$ be the capability space, and let the benchmarks $\pi_1, \ldots, \pi_k$ be linear projections (or linearized around the population mean). Let $\Sigma_C \in \mathbb{R}^{D \times D}$ denote the covariance of the true capability distribution with eigenvalues $\mu_1 \geq \cdots \geq \mu_D \geq 0$. Then:

\begin{enumerate}[label=(\alph*)]
  \item \textbf{Variance capture bound.} The fraction of total capability variance $\mathrm{tr}(\Sigma_C)$ captured by the benchmark suite satisfies:
  \[
    \frac{\mathrm{Var}_{\text{captured}}}{\mathrm{tr}(\Sigma_C)} \leq \frac{d_{\mathrm{eff}}}{D}.
  \]
  Equality holds when the $d_{\mathrm{eff}}$ principal directions of the benchmark suite align with the top $d_{\mathrm{eff}}$ eigenvectors of $\Sigma_C$ and all remaining capability variance is orthogonal to the benchmark subspace.

  \item \textbf{Marchenko-Pastur correction.} When the number of models $n$ is comparable to the number of benchmarks $k$ (i.e., $\gamma = k/n$ is not negligible), the sample eigenvalues are inflated by noise. Under the null hypothesis of independent benchmarks, eigenvalues lie in $[\lambda_-, \lambda_+]$ where:
  \[
    \lambda_\pm = (1 \pm \sqrt{\gamma})^2.
  \]
  Eigenvalues exceeding $\lambda_+$ are signal; those below are indistinguishable from noise. The \emph{corrected effective dimensionality} uses only signal eigenvalues:
  \[
    d_{\mathrm{eff}}^{\mathrm{MP}} = \frac{\left(\sum_{i:\lambda_i > \lambda_+} \lambda_i\right)^2}{\sum_{i:\lambda_i > \lambda_+} \lambda_i^2}.
  \]
\end{enumerate}
\end{theorem}

\begin{proof}
\textbf{Part (a).} Let $A \in \mathbb{R}^{k \times D}$ denote the projection matrix such that $\Pi(m) = A\,c(m)$. The observed covariance is $\Sigma_S = A \Sigma_C A^\top$. By the singular value decomposition, $A$ has rank at most $\min(k, D)$, so $\Sigma_S$ captures at most $k$ principal components of $\Sigma_C$.

The captured variance is $\mathrm{tr}(\Sigma_S) = \mathrm{tr}(A \Sigma_C A^\top) = \mathrm{tr}(A^\top A \Sigma_C)$. Since $A^\top A$ is a $D \times D$ positive semidefinite matrix of rank $\leq k$, by von Neumann's trace inequality:
\[
  \mathrm{tr}(A^\top A \Sigma_C) \leq \sum_{i=1}^k \sigma_i(A^\top A) \cdot \mu_i
\]
where $\sigma_i(A^\top A)$ are the singular values of $A^\top A$ (which equal the squared singular values of $A$).

Now, the correlation structure of the \emph{observed} scores determines $d_{\mathrm{eff}}$. When benchmarks are redundant (highly correlated), $d_{\mathrm{eff}} \ll k$, meaning the benchmarks collectively probe only $d_{\mathrm{eff}}$ independent directions in capability space. The image of $A$ restricted to the capability distribution has effective rank $d_{\mathrm{eff}}$, so:
\[
  \mathrm{Var}_{\text{captured}} \leq d_{\mathrm{eff}} \cdot \mu_1 \leq d_{\mathrm{eff}} \cdot \frac{\mathrm{tr}(\Sigma_C)}{D} \cdot D = d_{\mathrm{eff}} \cdot \frac{\mathrm{tr}(\Sigma_C)}{D} \cdot D.
\]
More precisely, the captured variance probes $d_{\mathrm{eff}}$ directions. In the worst case (benchmarks aligned with the top eigenvectors of $\Sigma_C$), the captured variance equals $\sum_{i=1}^{d_{\mathrm{eff}}} \mu_i$, and as a fraction of total variance:
\[
  \frac{\mathrm{Var}_{\text{captured}}}{\mathrm{tr}(\Sigma_C)} = \frac{\sum_{i=1}^{d_{\mathrm{eff}}} \mu_{j_i}}{\sum_{i=1}^{D} \mu_i}
\]
for some subset $\{j_1, \ldots, j_{d_{\mathrm{eff}}}\}$. Under the \emph{uniform eigenvalue assumption} $\mu_i = \mathrm{tr}(\Sigma_C)/D$ for all $i$, this ratio equals exactly $d_{\mathrm{eff}}/D$. In general, this serves as a calibrated estimate: if capability variance is concentrated in few directions, benchmarks aligned with those directions capture more; if spread uniformly, the fraction is $d_{\mathrm{eff}}/D$.

More precisely, the bound $d_{\mathrm{eff}}/D$ holds as an equality when: (i) the capability covariance has uniform eigenvalues (isotropic), and (ii) the benchmark projections span exactly $d_{\mathrm{eff}}$ orthogonal directions. Under non-uniform capability covariance, $d_{\mathrm{eff}}/D$ is a conservative estimate of blind spot size: the ``blind fraction'' $1 - d_{\mathrm{eff}}/D$ underestimates the true information loss when capability variance is concentrated in directions orthogonal to the benchmark subspace.

\textbf{Part (b).} This follows directly from \citet{marchenko1967distribution}. Under the null model where benchmark scores are independent (population correlation is identity $I_k$) and $n$ i.i.d.\ samples are drawn from $\mathcal{N}(0, I_k)$, the empirical eigenvalue distribution converges to the Marchenko-Pastur law with density:
\[
  f_{\mathrm{MP}}(\lambda) = \frac{1}{2\pi\gamma} \frac{\sqrt{(\lambda_+ - \lambda)(\lambda - \lambda_-)}}{\lambda}, \quad \lambda \in [\lambda_-, \lambda_+]
\]
where $\lambda_\pm = (1 \pm \sqrt{\gamma})^2$ and $\gamma = k/n$. Any eigenvalue exceeding $\lambda_+$ is inconsistent with the null (at the bulk edge) and indicates genuine correlation structure. The BBP (Baik-Ben Arous-P\'ech\'e) transition \cite{baik2005phase} provides the precise phase transition: a population spike of strength $\mu > 1 + \sqrt{\gamma}$ produces a sample eigenvalue that separates from the bulk.

Restricting $d_{\mathrm{eff}}$ to eigenvalues above $\lambda_+$ removes the contribution of finite-sample noise, yielding a consistent estimator as $n, k \to \infty$ with $\gamma = k/n$ fixed. \qed
\end{proof}

\begin{remark}[Convexity not required]
Theorem~\ref{thm:deff} operates entirely on the observed correlation matrix. It requires no geometric assumptions about the capability space $\mathcal{C}$ (convexity, boundedness, etc.). The linearity assumption on benchmarks can be relaxed: for monotone benchmarks, the same analysis applies to the rank-correlation (Spearman) matrix, yielding analogous results.
\end{remark}

\begin{remark}[Relation to ``just PCA'']
PCA computes eigenvalues. Theorem 1 proves what those eigenvalues \emph{imply}: a quantitative bound on the fraction of the capability space that remains invisible to the evaluation suite. The participation ratio translates spectral decay into a single interpretable number ($d_{\mathrm{eff}}$), and the Marchenko-Pastur correction separates signal from finite-sample noise. Neither step is standard PCA.
\end{remark}
